\documentclass[output=paper,colorlinks,citecolor=brown]{langscibook}
\ChapterDOI{10.5281/zenodo.21237009}
\author{Romain Garnier\orcid{}\affiliation{Université de Limoges}}
%\ORCIDs{}

\title{Delocatival adjectives surfacing as participles}

\abstract{This paper proposes to reassess several active participles of the type *\textit{CC-en-i̯ónt}- and \textit{*CC-ént-} as decasuatives built on a “petrified” locative or oblique *\textit{CC-én}. At its face value, Ved.\ \textit{bhuraṇ-yánt}- ‘quivering, active’ (of Agni) is the participle of the intransitive present stem \textit{bhuraṇ-yá-\textsuperscript{ti}} ‘to be active or restless’, alongside its deverbative doublet \textit{bhuraṇ-yú}- [adj.] ‘quivering, active’. In fact, the indicative \textit{bhuraṇ-yá-\textsuperscript{ti}} is likely to have been backformed  after \textit{bhuraṇ-yánt}-, which may be accounted for as a decasuative adjective (reflecting PIE *\textit{bʰr̥h\textsubscript{2}-en-i̯ó-nt}-), derived from a locative *\textit{bʰr̥h\textsubscript{2}-én} ‘in quick motion’. The suffix *-\textit{nt}- was added as a characterizing morpheme to a former *-\textit{i̯ó}-stem, with the same pattern as PIE *\textit{g̑erh\textsubscript{2}-ó-nt-es} ‘the old ones’ (vs.\ simplex *\textit{g̑erh\textsubscript{2}-ó-}) or *\textit{g\textsuperscript{u̯}ih\textsubscript{3}-u̯ó-nt-es} ‘the living ones’ (vs.\ simplex *\textit{g\textsuperscript{u̯}ih\textsubscript{3}-u̯ó-}). Several other cases will be dealt with, involving PIE 1.\ *\textit{bʰeh\textsubscript{2}}- ‘to shine’, PIE 2.\ *\textit{bʰeh\textsubscript{2}}- ‘to speak, to say’, and PIE *\textit{g\textsuperscript{u̯}eh\textsubscript{2}}- ‘to make a step’.}


\IfFileExists{../localcommands.tex}{
   \addbibresource{../localbibliography.bib}
   \usepackage{tabularx,multicol}
\usepackage{url}
\urlstyle{same}

 
\usepackage{langsci-optional}
\usepackage{langsci-lgr}
\usepackage{langsci-gb4e}
\usepackage[linguistics,edges]{forest}
\usepackage{tikz-qtree}
\usetikzlibrary{positioning}
\usetikzlibrary{patterns.meta}
\usepackage{multirow}
\usepackage{pgfplots}
\usepackage{setspace}
\usepackage{amsmath,stackengine}
% \usepackage{fontspec}
% \usepackage{tipa} % TIPA is banned
\usepackage{langsci-textipa}
\usepackage{langsci-branding}
\usepackage{longtable}
\usepackage{tabto}

\usepackage{enumitem}

%\usepackage{lscape} %Querformat chapter 4

   % \newcommand*{\orcid}{}


\makeatletter
\let\thetitle\@title
\let\theauthor\@author
\makeatother

\newcommand{\togglepaper}[1][0]{
   \bibliography{../localbibliography}
   \papernote{\scriptsize\normalfont
     \theauthor.
     \titleTemp.
     To appear in:
    Laura Grestenberger, Viktoria Reiter \& Melanie Malzahn (eds.).
    \emph{Deadjectival verb formation in Indo-European and beyond}.
     Berlin: Language Science Press. [preliminary page numbering]
   }
   \pagenumbering{roman}
   \setcounter{chapter}{#1}
   \addtocounter{chapter}{-1}
}

\newbool{bookcompile}
\booltrue{bookcompile}
\newcommand{\bookorchapter}[2]{\ifbool{bookcompile}{#1}{#2}}

%Sonderzeichen
%idg palatales k
\newcommand{\kinvbreve}{k\raisebox{0.6ex}{\hspace{-0.05em}\char"0311}}

\newcommand{\jac}{%
  \textit{\char"0237}% italic dotless j (ȷ)
  \hspace{-0.01em}% fine-tune horizontal spacing
  \raisebox{0.01ex}{\char"0301}% raised combining acute
}

%Baltic circumflex l
\newcommand\ltilded{%
\hspace{-0.2em}%
  \stackon[-1.5pt]{l}{\smash{\kern0.2em\~{}}}%
  \hspace{-0.2em}%
}

%Slav. wedge+ double acute
\newcommand{\HighH}[1]{%
    \stackon[-1.2ex]{#1}{\kern0.1em\H{}\kern-0.1em}%
}

%Lith dot tilde VICKY diese b.-sl. Sonderzeichen BRINGEN MICH UM!!
\newcommand{\tildedot}[1]{%
    \ensurestackMath{%
        \stackon[-0.27ex]{% Tilde layer 
            \stackon[-0.15ex]{% Dot layer 
                #1%
            }{\hspace{0.1em}\cdot\hspace{-0.13em}}%
        }{\hspace{0.1em}\text{\~{}}\hspace{-0.13em}%
    }%
}}

%Lith dot acute
\newcommand{\dotacute}[1]{%
    \ensurestackMath{%
        \stackon[-1.1ex]{% acute layer 
            \stackon[-0.15ex]{% Dot layer 
                #1%
            }{\hspace{0.1em}\cdot\hspace{-0.12em}}%
        }{\hspace{0.1em}\text{\'{}}\hspace{-0.12em}%
    }%
}}



% \setlength{\emergencystretch}{2pt} 

 \pgfplotsset{compat=1.18}


\defcitealias{AbB_1}{AbB 1}
\defcitealias{AbB_7}{AbB 7}
\defcitealias{ABL_1}{ABL 1}
\defcitealias{Vries1977}{AEW}
\defcitealias{ALEW}{ALEW}
\defcitealias{AnGr}{AnGr}
\defcitealias{ARM_1}{ARM 1}
\defcitealias{BT}{BT}
\defcitealias{black_concise_2000}{CDA}
\defcitealias{CHD}{CHD}
\defcitealias{CT28}{CT 28}
\defcitealias{CT34}{CT 34}
\defcitealias{CT39}{CT 39}
\defcitealias{Cleasby1874}{CV}
\defcitealias{Beekes2010}{EDG}
\defcitealias{eDiAna}{eDiAna}
\defcitealias{eDIL}{eDIL}
\defcitealias{EDL}{EDL}
\defcitealias{Kroonen2013}{EDPG}
\defcitealias{ESJS}{ESJS}
\defcitealias{EVP}{EVP}
\defcitealias{EWA1988}{EWA}
\defcitealias{EWAI}{EWAia I}
\defcitealias{EWAII}{EWAia II}
\defcitealias{EWGP}{EWGP}
\defcitealias{GDLI1961}{GDLI}
\defcitealias{GPC}{GPC}
\defcitealias{GRADIT1999}{GRADIT}
\defcitealias{HW}{HW}
\defcitealias{IEW}{IEW}
\defcitealias{KAH_2}{KAH 2}
\defcitealias{KAR_1}{KAR 1}
\defcitealias{KAR_2}{KAR 2}
\defcitealias{KEWA}{KEWA}
\defcitealias{Labat_TDP}{Labat TDP}
\defcitealias{lambert_BWL}{Lambert BWL}
\defcitealias{LEIA}{LEIA}
\defcitealias{Rix2001}{LIV\textsuperscript{2}}
\defcitealias{LIVaddenda}{LIV\textsuperscript{2} {A}ddenda}
\defcitealias{LKG1971}{LKG}
\defcitealias{LKZ}{LKŽ}
\defcitealias{LP}{LP}
\defcitealias{MhdGr}{MhdGr}
\defcitealias{MLLVG1959}{MLLVG}
\defcitealias{NIL}{NIL}
\defcitealias{OgnS}{OgnS}
\defcitealias{ONP}{ONP}
\defcitealias{RIMA_2}{RIMA 2}
\defcitealias{TIM_2}{TIM 2}
\defcitealias{WAP}{WAP}
\defcitealias{YOS_2}{YOS 2}






\newindex{wan}{wdx}{wnd}{Form index}
\newcommand{\iw}[2]{\index[wan]{#1!{#2}\xspace}}
\newcommand{\iwi}[3][]{\index[wan]{#2!{#3}\xspace#1}{#3}}
\newcommand{\iwit}[3][]{\rule{0pt}{0pt}#1\index[wan]{#2!{#3}}{#3}}


\newindex{van}{vdx}{vnd}{Form index}
\newindex{ban}{bdx}{bnd}{Book index}

% \newcommand{\iw}[3][]{\index[wan]{#2!\textit{#3}}}


% \providecommand{\iwi}[3][]{\iw{#3}\textit{#3}}


\newcommand{\indexnote}[1]{{\footnotesize\sffamily\raggedright\normalfont\color{gray}#1}}

   %% hyphenation points for line breaks
%% Normally, automatic hyphenation in LaTeX is very good
%% If a word is mis-hyphenated, add it to this file
%%
%% add information to TeX file before \begin{document} with:
%% %% hyphenation points for line breaks
%% Normally, automatic hyphenation in LaTeX is very good
%% If a word is mis-hyphenated, add it to this file
%%
%% add information to TeX file before \begin{document} with:
%% %% hyphenation points for line breaks
%% Normally, automatic hyphenation in LaTeX is very good
%% If a word is mis-hyphenated, add it to this file
%%
%% add information to TeX file before \begin{document} with:
%% \include{localhyphenation}
\hyphenation{
    par-a-digm non-de-adj-ect-iv-al -pre-sent -pre-sents for-ma-tion-al in-do-ger-ma-ni-sche in-do-ger-ma-ni-schen Rei-chert Rei-ter dia-chrony ety-mo-lo-gi-sches Ost-ger-ma-ni-schen ger-ma-ni-schen alt-in-di-schen mor-pho-lo-gi-scher Alt-indo-ari-sch-en lexico-graphy Schrij-ver
}
% aktuell in Bibliographie: stud-ies, histor-ical Soci-età
\hyphenation{
    par-a-digm non-de-adj-ect-iv-al -pre-sent -pre-sents for-ma-tion-al in-do-ger-ma-ni-sche in-do-ger-ma-ni-schen Rei-chert Rei-ter dia-chrony ety-mo-lo-gi-sches Ost-ger-ma-ni-schen ger-ma-ni-schen alt-in-di-schen mor-pho-lo-gi-scher Alt-indo-ari-sch-en lexico-graphy Schrij-ver
}
% aktuell in Bibliographie: stud-ies, histor-ical Soci-età
\hyphenation{
    par-a-digm non-de-adj-ect-iv-al -pre-sent -pre-sents for-ma-tion-al in-do-ger-ma-ni-sche in-do-ger-ma-ni-schen Rei-chert Rei-ter dia-chrony ety-mo-lo-gi-sches Ost-ger-ma-ni-schen ger-ma-ni-schen alt-in-di-schen mor-pho-lo-gi-scher Alt-indo-ari-sch-en lexico-graphy Schrij-ver
}
% aktuell in Bibliographie: stud-ies, histor-ical Soci-età
   \boolfalse{bookcompile}
   \togglepaper[5]%%chapternumber
}{}

\begin{document}
\SetupAffiliations{mark style=none}
\maketitle


\section{Vedic -\textit{anyá}-presents: state of the art}
\label{sec:sota}

In Vedic, there is a small class of -\textit{anyá}-presents \iw{Sanskrit}{-\textit{anyá}-} such as \textit{tur-aṇ-yá}-\textit{\textsuperscript{ti}} \iw{Sanskrit}{\textit{turaṇyá}-} ‘to rush’, \textit{bhur-aṇ-yá-\textsuperscript{ti}} \iw{Sanskrit}{\textit{bhuraṇyá}-} ‘to be active or restless’ (\isi{transitive} ‘to agitate, to stir [a liquid]’), \textit{riṣ-aṇ-yá}-\textit{\textsuperscript{ti}} \iw{Sanskrit}{\textit{riṣaṇyá}-} ‘to suffer injury’, \textit{kr̥p-an-yá-\textsuperscript{ti}} \iw{Sanskrit}{\textit{kr̥panyá}-} `to desire’, \textit{pr̥t-an-yá-\textsuperscript{ti}} \iw{Sanskrit}{\textit{pr̥tanyá}-} ‘to fight’. YAv.\ \textit{zar-an-iia}- \iw{Avestan}{\textit{zaraniia}-} ‘to be irritated’ reflects the Iranian counterpart of this verbal suffix. One may note that most of these verbs are \isi{intransitive} (Ved.\ \textit{bhur-aṇ-yá-\textsuperscript{ti}} \iw{Sanskrit}{\textit{bhuraṇyá}-} only being labile).\is{labile verbs} \citet[122]{Jasanoff2003} suggests identifying this verbal derivation\is{derivation!verbal} with Hittite \isi{durative}s in \iwi{Hitt}{-\textit{anna/i}-}, by assuming that the geminates reflect PIE \iwi{PIE}{*-\textit{nh\textsubscript{2}-i}-}. To trigger the assimilation *-\textit{nH}- > -\textit{nn}-, the suffix requires an initial vowel, reconstructed as PIE \iwi{PIE}{*-\textit{énh\textsubscript{2}-i}-} by \citet[40--43]{Melchert1992} or \iwi{PIE}{*-\textit{ótn-i}-} by \citet[175]{Kloekhorst2008}. According to \citet[125]{Jasanoff2003}, Ved.\ -\textit{anyá}-presents \iw{Sanskrit}{-\textit{anyá}-} cannot be denominatives\is{denominative!verb} to nominal formations in \iwi{Sanskrit}{-\textit{an}-} (or \iwi{Sanskrit}{-\textit{ana}-}) because the underlying nasal stems are never attested – a claim which needs to be reassessed. Besides – as he duly points out – thematic nouns in \iwi{Sanskrit}{-\textit{ana}-} would rather produce denominatives\is{denominative!verb} in **-\textit{anā̆-yá-}\textit{\textsuperscript{ti}}. Jasanoff instead compares Ved.\ -\textit{anyá}-presents \iw{Sanskrit}{-\textit{anyá}-} to the type of Ved.\ \textit{gr̥bh-ā-yá-\textsuperscript{ti}} \iw{Sanskrit}{\textit{gr̥bhāyá}-} ‘to seize’ (< PIE *-\textit{n̥h\textsubscript{2}-i̯é/ó}-),\iw{PIE}{*-\textit{n̥h\textsubscript{2}-i̯e/o}-} by assuming an unrecognized Indo-Iranian sound change triggering the shortening of a long vocalic sonorant when following another vocalic sonorant (\citealt[126]{Jasanoff2003}). He claims that Ved.\ \textit{dam-an-yá-\textsuperscript{ti}}\iw{Sanskrit}{\textit{damanyá}-} ‘to subdue’ would be the regular reflex of PIE \iwi{PIE}{*\textit{dm̥h\textsubscript{2}-n̥h\textsubscript{2}-i̯é/ó}-}, whereas the \isi{doublet} \textit{dam-ā-yá-\textsuperscript{ti}} \iw{Sanskrit}{\textit{damāyá}-} ‘id.’ would be analogically remade after the pattern of \textit{gr̥bh-ā-yá-\textsuperscript{ti}} \iw{Sanskrit}{\textit{gr̥bhāyá}-} ‘to seize’. This seems to me an \textit{ad hoc} rule, which is not supported by any independent lexeme. Furthermore, this implies that forms such as Ved.\ \textit{riṣ-aṇ-yá}-\textit{\textsuperscript{ti}} \iw{Sanskrit}{\textit{riṣaṇyá}-} ‘to suffer injury’, \textit{kr̥p-an-yá-\textsuperscript{ti}} \iw{Sanskrit}{\textit{kr̥panyá}-} ‘to desire’, or \textit{pr̥t-an-yá-\textsuperscript{ti}} \iw{Sanskrit}{\textit{pr̥tanyá}-} ‘to fight’ would be analogical. This phonological account should be rejected, since Jasanoff does not explain why Vedic -\textit{anyá}-presents \iw{Sanskrit}{-\textit{anyá}-} (always inflected in the active) are often \isi{intransitive} or labile.\is{labile verbs} By contrast, -\textit{n{ā́}-}\textit{\textsuperscript{ti}} \iw{Sanskrit}{-\textit{n{ā́}}-} or -\textit{ā-yá-\textsuperscript{ti}} \iw{Sanskrit}{-\textit{āyá}-} presents are exclusively \isi{transitive} (for instance Ved.\ \textit{iṣ-ṇ{ā́}-}\textit{\textsuperscript{ti}} \iw{Sanskrit}{\textit{iṣṇ{ā́}}-} ‘to send’). The alleged lack of evidence for -\textit{an}-stems \iw{Sanskrit}{-\textit{an}-} does not withstand scrutiny: alongside \textit{iṣ-aṇ-yá}-\textit{\textsuperscript{ti}} \iw{Sanskrit}{\textit{iṣaṇyá}-} ‘to cause to make haste, excite, drive’, there is \textit{iṣ-áṇi} \iw{Sanskrit}{\textit{iṣáṇi} (loc.)} [loc.] (RV 2.2.9d) ‘at [her] impulsion’. Likewise, alongside \textit{tur-aṇ-yá}-\textit{\textsuperscript{ti}} \iw{Sanskrit}{\textit{turaṇyá}-} ‘to be quick or swift’, there is \textit{tur-áṇa-} \iw{Sanskrit}{\textit{turáṇa}-} [adj.] ‘swift’ and \textit{bhur-aṇa}- \iw{Sanskrit}{\textit{bhuraṇa}-} [adj.] ‘quick, active’ (= stressed \iw{Sanskrit}{*\textit{bhuráṇa}-} *\textit{bhur-áṇa}-), akin to \textit{bhur-aṇ-yá-\textsuperscript{ti}} \iw{Sanskrit}{\textit{bhuraṇyá}-} ‘to be active or restless’ and \textit{bhur-aṇ-yú}- \iw{Sanskrit}{\textit{bhuraṇyú}-} [adj.] ‘quivering, active’.

\largerpage
Independently, the comparison of Ved.\ -\textit{anyá}-presents \iw{Sanskrit}{-\textit{anyá}-} with Hitt.\ \iwi{Hitt}{-\textit{anni}-} must now be considered outdated, since this Anatolian \isi{iterative} suffix cannot reflect PIE \iwi{PIE}{*-\textit{n̥h\textsubscript{2}-i}-} (or \iwi{PIE}{*-\textit{énh\textsubscript{2}-i}-}), but rather PA *-\textit{{ā́}d-n-i}-, \iw{PAn}{*-\textit{{ā́}d-n-i}-} as \textit{per} \citet[526--528]{Sasseville2021}. In my opinion, one has to start with an archaic endingless locative such as Ved.\ \iwi{Sanskrit}{*\textit{iṣ-án}} ‘with drive, impetus’,  \iwi{Sanskrit}{*\textit{bhur-án}} ‘in quick motion’, Av.\  \iwi{Avestan}{*\textit{zar-an}} ‘with anger’. The lability\is{labile verbs} of the derived verbs vindicates their nominal origin (viz.\ \textit{bhur-aṇ-yá-\textsuperscript{ti}} \iw{Sanskrit}{\textit{bhuraṇyá}-} ‘to be/to set in quick motion’), since  \iwi{Sanskrit}{*\textit{iṣ-án}} ‘with drive, impetus’ could be either agent- or patient-oriented. As a result, several “active” participles\is{participle} of the type *\textit{CC-en-i̯ónt}-  \iw{PIE}{*-\textit{en-i̯ó-nt}-} could in fact be reanalyzed as old decasuative adjectives\is{decasuative} built on a “petrified” locative or oblique stem \iw{PIE}{*-\textit{én}} *\textit{CC-én}. Another strategy of coining a delocative adjective\is{delocatival} was to add a thematic vowel: PIE \iwi{PIE}{*-\textit{én-o}-} (cf.\ Ved.\ \textit{tur-áṇ-a-} \iw{Sanskrit}{\textit{turáṇa}-} [adj.] ‘swift’ and *\textit{bhur-áṇ-a}- \iw{Sanskrit}{*\textit{bhuráṇa}-} [adj.] ‘quick, active’ < PIE \iwi{PIE}{*\textit{bʰr̥h\textsubscript{2}-én-o}-}). I would tentatively account for Ved.\ \textit{dam-an-yá-\textsuperscript{ti}} \iw{Sanskrit}{\textit{damanyá}-} ‘to subdue’ by reconstructing a PIE \iwi{PIE}{*\textit{démh\textsubscript{2}-mn̥}} [nt.] ‘taming’, which was prone to undergo the so-called “\textit{ašnō}-rule”, \is{\textit{ašnō}-rule} viz.\ the reduction of PIE \iwi{PIE}{*-\textit{mn}-} to *\textit{-n}- after obstruent, cf.\ \citet{Nussbaum2010} and \citet[186--188]{Melchert2010}. The \isi{genitive} singular *\textit{démh\textsubscript{2}-(m)n-os} would thus surface as Ved.\ \iwi{Sanskrit}{*\textit{dám-n-as}} after which an endingless locative \iwi[(loc.)]{Sanskrit}{*\textit{dám-an}} ‘in subduing’ was formed. This very form may have been the base of the denominative verb\is{denominative!verb} \textit{dam-an-yá-\textsuperscript{ti}} \iw{Sanskrit}{\textit{damanyá}-} ‘to subdue’ (with denominal \is{denominal!verb} \is{causative}  causativity, viz.\ *‘to put into subduing, to cause sbdy.\ or sth.\ to be in the state of being subdued’). One may reasonably speculate that \textit{dam-an-yá-\textsuperscript{ti}} \iw{Sanskrit}{\textit{damanyá}-} ‘to subdue’ is basically an indicative and was built within Vedic, by contrast with \textit{bhur-aṇ-yá-\textsuperscript{ti}}, \iw{Sanskrit}{\textit{bhuraṇyá}-} which is likely to have been inherited, since there is comparative evidence for PIE \iwi{PIE}{*\textit{bʰr̥h\textsubscript{2}-én}} [loc.] ‘in quick motion’ (see Section \ref{sec:2}).

As already mentioned, Ved.\ \textit{bhur-aṇ-yú}- \iw{Sanskrit}{\textit{bhuraṇyú}-} [adj.] ‘quivering, active’ is the functional equivalent of the \isi{participle} \textit{bhur-aṇ-yánt}- \iw{Sanskrit}{\textit{bhuraṇyánt}-} ‘id.’, like \textit{tur-aṇ-yú}- \iw{Sanskrit}{\textit{turaṇyú}-} [adj.] ‘rushing’ vs.\ \textit{tur-aṇ-yánt}- \iw{Sanskrit}{\textit{turaṇyánt}-} ‘id.’. One has perhaps to assume that the \isi{intransitive} or labile\is{labile verbs} -\textit{anyá}-presents \iw{Sanskrit}{-\textit{anyá}-} were backformed \is{backformation} after their participles\is{participle} \iwi{Sanskrit}{-\textit{an-yánt}-}, which were used as adjectives (‘rushing, quivering’). The starting point may have been an endingless locative \iwi{Sanskrit}{*\textit{bhur-án}} ‘in quick motion’ (< PIE \iwi{PIE}{*\textit{bʰr̥h\textsubscript{2}-én}}) from which was derived (i) a decasuative\is{decasuative} \iwi{Sanskrit}{-\textit{yánt}-}\isi{participle}, and (ii) a decasuative \iwi{Sanskrit}{-\textit{yú}-}adjective.\is{decasuative} I consider that in such cases the \iwi{Sanskrit}{-\textit{yánt}-}\isi{participle} is to be segmented as *-\textit{yá}-\textit{nt}-, with the well-known “characterizing” morpheme\iw{PIE}{*-\textit{(e/o)nt}-} *-\textit{nt}-\footnote{See for instance PIE *\textit{g̑erh\textsubscript{2}-ó-nt-es} \iw{PIE}{*\textit{g̑erh\textsubscript{2}-ó-nt}-} ‘the elders, the old ones’ (vs.\ simplex \iwi{PIE}{*\textit{g̑erh\textsubscript{2}-ó}-} [adj.] ‘wrinkly, decayed’) or PIE *\textit{g\textsuperscript{u̯}ih\textsubscript{3}-u̯ó-nt-es} \iw{PIE}{*\textit{g\textsuperscript{u̯}ih\textsubscript{3}-u̯ó-nt}-} ‘the living ones, the living world’ (vs.\ simplex \iwi{PIE}{*\textit{g\textsuperscript{u̯}ih\textsubscript{3}-u̯ó}-} [adj.] ‘quick, alive’). The indicative \iw{PIE}{*\textit{g\textsuperscript{u̯}ih\textsubscript{3}-u̯é/ó}-} *\textit{g\textsuperscript{u̯}ih\textsubscript{3}-u̯-é/ó}- ‘to live’ (Ved.\ \iw{Sanskrit}{\textit{jī́va}-} \textit{jī́va-\textsuperscript{ti}}, Lat.\ \textit{vīvō}\iw{Latin}{\textit{uīuō}}) is likely to have been backformed \is{backformation} after its (pseudo-)\isi{participle} stem. That the\iw{PIE}{*-\textit{(e/o)nt}-} *-\textit{ont}-participles\is{participle} were formerly independent adjectival stems is assumed by \citet{Melchert2017}. For the suffix, one may compare CLuw.\ \iwi{CLuw}{\textit{wayant(i)}-} [c.] ‘animal’ from PIE  *\textit{g\textsuperscript{u̯}oi̯h\textsubscript{3}-ó-nt-es} \iw{PIE}{*\textit{g\textsuperscript{u̯}oi̯h\textsubscript{3}-ó-nt}-} ‘the living ones; creatures’ which is a recharacterization of PIE \iwi{PIE}{*\textit{g\textsuperscript{u̯}oi̯h\textsubscript{3}-ó}-}  [adj.] ‘living’. PIE *\textit{g\textsuperscript{u̯}ih\textsubscript{3}-u̯ó-nt-es} \iw{PIE}{*\textit{g\textsuperscript{u̯}ih\textsubscript{3}-u̯ó-nt}-} is the \textit{Scharnierform} yielding a present \iw{PIE}{*\textit{g\textsuperscript{u̯}ih\textsubscript{3}-u̯é/ó}-} *\textit{g\textsuperscript{u̯}ih\textsubscript{3}-u̯-é/ó}-, which can hardly be accounted for by \isi{conversion} \iwi{PIE}{*\textit{g\textsuperscript{u̯}ih\textsubscript{3}-u̯ó}-} → *\textit{g\textsuperscript{u̯}íh\textsubscript{3}-u̯-e/o}- or by deletion (viz.\ PIE *\textit{g\textsuperscript{u̯}íh\textsubscript{3}-u̯(e/o-)e/o-}), \textit{pace} \citet[273--274]{VillanuevaSvensson2021}.} added to a decasuative\is{decasuative} \mbox{*-\textit{yá}-}adjective\iw{Sanskrit}{*-\textit{yá}- (adj.)} (< PIE \iwi{PIE}{*-\textit{i̯ó}-}). To sum up, one has to start with PIE \iwi{PIE}{*\textit{bʰr̥h\textsubscript{2}-en-i̯ó}-} [adj.]  ‘quivering’,\footnote{The same pattern is seen in PIE \iwi{PIE}{*\textit{g̑ʰl̥h\textsubscript{3}-en-i̯ó}-} [color adj.] ‘yellow’ → Ved.\ \iwi{Sanskrit}{\textit{híraṇya}-} [nt.] ‘gold’ (lit.\ ‘being in the state of shining’).} a decasuative stem\is{decasuative} which was recharacterized as PIE \iwi{PIE}{*\textit{bʰr̥h\textsubscript{2}-en-i̯ó-nt}-} ‘the swift one’ (Ved.\ \textit{bhur-aṇ-yánt}- \iw{Sanskrit}{\textit{bhuraṇyánt}-} ‘quivering, active’ [of Agni]). Independently, there was also PIE \iwi{PIE}{*\textit{bʰr̥h\textsubscript{2}-en-i̯ú}-} [adj.] (Ved.\ \iw{Sanskrit}{\textit{bhuraṇyú}-} \textit{bhur-aṇ-yú}-) and also PIE \iwi{PIE}{*\textit{bʰr̥h\textsubscript{2}-én-o}-} [adj.] (Ved.\ \iwi{Sanskrit}{*\textit{bhuráṇa}-}).

\largerpage
The derivational history of Ved.\ \textit{bhur-aṇ-yá-\textsuperscript{ti}} \iw{Sanskrit}{\textit{bhuraṇyá}-} proposed above is not totally unparalleled: There are good Anatolian parallels. Hitt.\ *\textit{išḫan-iya}- \iw{Hitt}{*\textit{išḫaniya}-} ‘to bind’ reconstructed by \citet[286--287]{Rieken1999} based on Hitt.\ \iwi{Hitt}{\textit{išḫanittar}} [c.] ‘relative by marriage’ (< PA \iwi{PAn}{*\textit{sH-an-y-é-tros}}) requires a base \iwi{PAn}{*\textit{išḫan}} (< PIE \iwi{PIE}{*\textit{sh\textsubscript{2}-én}} [loc./obl.] ‘with bond’), which she further connects with Hitt.\ \iwi{Hitt}{\textit{šaḫḫan}-} [nt.] ‘(feudal) obligation’.\footnote{One should note that, from now on, due to the existence of the so-called \is{\textit{ašnō}-rule} “\textit{ašnō}-rule”, Hitt.\ \iwi{Hitt}{\textit{šaḫḫan}-} shall no longer be directly assigned to PIE \textsuperscript{†}\textit{séh\textsubscript{2}}\textit{-}{\textit{n̥}} ‘obligation’ \textit{pace} \citet[287]{Rieken1999} and \citet[691]{Kloekhorst2008}, but to PIE \iwi{PIE}{*\textit{séh\textsubscript{2}-mn̥}} [nt.] with nasal-deletion after the \is{genitive}gen.sg.\ *\textit{séh\textsubscript{2}-(m)n-os}.}  Interestingly, “Luwoid” derivatives in \iwi{Hitt}{-\textit{att-allā}-} (with \iwi{Hitt}{-\textit{allā}-} from PA *\mbox{-\textit{élā}-)} \iw{PAn}{*-\textit{élā}-} and \iwi{Hitt}{-\textit{alla/i}-} (< PA \iwi{PAn}{*-\textit{élo}-} [with the so-called “Luwian \textit{i}-mutation”]) \is{\textit{i}-mutation} are  built directly on the stem of the base: Hitt.\ \iwi{Hitt}{\textit{išḫan-attallā}-} [c.] (?) and \iwi{Hitt}{\textit{išḫan-alla/i}-} [c.] (?), both of unclear meaning. As I have recently suggested \citep[175]{Garnier2022b}, Hitt.\ \textit{kuššan-iye/a}-\textit{\textsuperscript{zzi}} \iw{Hitt}{\textit{kuššan-iye/a}-} ‘to hire, to employ’, which is akin to Proto-Gmc.\  *\textit{χūz-jan}- \iw{PGmc}{*\textit{χūzjan}-}  ‘to hire’ (\citealt[498]{Kloekhorst2008}; cf.\ OE \iwi{OE}{\textit{hȳr}} [str.f.] ‘wages, interest, usury’ < \iw{PGmc}{*\textit{χūzṓ}-}  *\textit{χūz-ṓ}-), may be accounted for by a base  \iwi{PAn}{*\textit{kuššan}} [loc.] ‘into [temporary] possession’, reflecting PIE \iwi{PIE}{*\textit{\kinvbreve uh\textsubscript{1}-s-én}} which was a \isi{doublet} of PIE \iwi{PIE}{*\textit{\kinvbreve uh\textsubscript{1}-és(i)}} [loc.], possibly associated to PIE \iwi{PIE}{*\textit{\kinvbreve éu̯h\textsubscript{1}-ōs}} [f.] ‘possession’ (of holokinetic inflection), or perhaps PIE \iwi{PIE}{*\textit{\kinvbreve éu̯h\textsubscript{1}-s}-} [nt.] ‘id.’ (of \isi{proterokinetic} inflection). The noun \iwi{Hitt}{\textit{kuššan}-} [nt.] ‘pay, salary, fee, hire’ is likely to have been backformed\is{backformation} after \textit{kuššan-iye/a}-\textit{\textsuperscript{zzi}}  \iw{Hitt}{\textit{kuššan-iye/a}-} ‘to hire, to employ’. For the meaning, one may compare Gk.\ \iwi{AG}{κῡριεία} [f.] ‘possession’ vs.\ \iwi{AG}{κῡριoς}  [m.] ‘lord, ruler’ (< PIE \iwi{PIE}{*\textit{\kinvbreve uh\textsubscript{1}-r-i-i̯ó}-} [adj.] ‘endowed with strength’, cf.\ PIE \iwi{PIE}{*\textit{\kinvbreve uh\textsubscript{1}-r-í}-} [f.] ‘strength’ derived from PIE \iwi{PIE}{*\textit{\kinvbreve uh\textsubscript{1}-r-ó}-} [adj.] ‘strong’). Interestingly, the “locative-1” (PA \iwi{PAn}{*\textit{\kinvbreve uʔ-és}}) is indirectly reflected by Lyd.\ \iwi{Lydian}{\textit{qašaas}} [c.] ‘hire; rental’ (< PA \iwi{PAn}{*\textit{\kinvbreve uʔ-es-y{ā́}}-}) and Lyc.\ \textit{qehñni-te} \iw{Lycian}{\textit{qehñni}-} ‘he rented’ (< Pre-Lyc.\ *\textit{qes-Vni-\textsuperscript{ti}}), with unlenited endings pointing to PA \iwi{PAn}{*-\textit{yé/ó}-} (\citealt[166--167]{Sasseville2021}). 

\section{A case study: PIE *\textit{bʰr̥h\textsubscript{2}}\textit{-én} ‘in quick motion’}\label{sec:2}
\subsection{PIE *\textit{bʰr̥h\textsubscript{2}}\textit{-én} ‘in quick motion’} 
\label{sec:2.1}

I have reconstructed a PIE etymon \iwi{PIE}{*\textit{bʰr̥h\textsubscript{2}-én}} [loc.] ‘in quick motion’ \citep[9]{Garnier2022a}, which would surface as Ved.\ *\textit{bhur-án} (for which see above). There is additional evidence for assuming such an “adverbial” stem: (i) PIE \iwi{PIE}{*\textit{bʰr̥h\textsubscript{2}-ḗn}} [HK f.] (< PIE *-\textit{én-s}) ‘vehemence’ (Section \ref{sec:2.2}); (ii) PIE \iwi{PIE}{*\textit{bʰr̥h\textsubscript{2}-én-o}-} [adj.] ‘being in quick motion’ (Section \ref{sec:2.3}); (iii) PIE \iwi{PIE}{*\textit{bʰr̥h\textsubscript{2}-én-t}-} [adj.] ‘being in quick motion’ (Section \ref{sec:2.4}); (iv) PIE \iwi{PIE}{*\textit{bʰr̥h\textsubscript{2}-en-i̯ó}-} [adj.] ‘quivering, active’ → “characterized” PIE \iwi{PIE}{*\textit{bʰr̥h\textsubscript{2}-en-i̯ó-nt}-} ‘the swift one’ and PIE \iwi{PIE}{*\textit{bʰr̥h\textsubscript{2}-en-i̯ú}-} [decasuative adj.]\is{decasuative}  ‘quivering, active’ (Section \ref{sec:2.5}). PIE *-\textit{én}-locatives \iw{PIE}{*-\textit{én}} from root nouns\is{root!noun} are assumed by \citet[326--327]{Neri2017}.

\subsection{PIE *\textit{bʰr̥h\textsubscript{2}}\textit{-ḗn} [HK f.] ‘rush, fury, vehemence’} 
\label{sec:2.2}

OIr.\ \iwi{OIr}{\textit{barae}} [f.] ‘fury, anger’ reflects Proto-Celt.\ \iwi{PCelt}{*\textit{bar-ēs}} (with *-\textit{ē} < *-\textit{en-s}), acc.sg.\ \iwi{PCelt}{*\textit{bar-en-an}},\footnote{Within Celtic a zero-graded secondary derivative\is{derivation!stem-based}  \iwi{PCelt}{*\textit{bar-n-iko}-} was coined, cf.\ OIr.\ \iwi{OIr}{\textit{barnech}/\textit{bairnech}} [adj.] ‘irritated’.}  which in turn reflects a remodeling of an old \isi{hysterokinetic} stem: PIE \iwi{PIE}{*\textit{bʰr̥h\textsubscript{2}-ḗn}} [f.] ‘vehemence’ (with *-\textit{ḗn} < PIE *-\textit{én-s}).\footnote{There was no loss of the final nasal after a long vowel due to the presence of the accent, as \textit{per} \citet[218--219]{Hardarson2005}. For the internal derivation yielding a \isi{hysterokinetic} stem from a PIE \iw{PIE}{*-\textit{én}} *-\textit{én}-locative, see \citet[324, fn.\ 181]{Neri2017}.} This inherited form (such nasal stems are not productive\is{productivity} in Celtic) has been assigned by Stüber \textit{apud} \citet[1--2]{Balles2002} to PIE \iwi{PIE}{*\textit{bʰerh\textsubscript{x}}-} ‘mit einem scharfen Werkzeug arbeiten’ (\citetalias[ 64]{Rix2001}), exemplified by Lat.\ \iwi{Latin}{\textit{feriō}} ‘to strike’ and \iwi{Latin}{\textit{forāre}} ‘to bore through, to pierce’. \citet[7]{Balles2002} further compares OIr.\ \iwi{OIr}{\textit{barae}} with Hom.\ \iwi{AG}{φρήν*} [f.] ‘midriff, diaphragm’ (in fact a \textit{plurale tantum} \iwi{AG}{φρένες}), which has almost nothing to do with ‘anger’.\footnote{I have explained Hom.\ \iwi{AG}{ἄφρων} [adj.] ‘senseless, foolish, heedless’  as the reflex of a PIE etymon \iwi{PIE}{*\textit{n̥-g\textsuperscript{u̯h}r(h\textsubscript{1})-on}-} ‘without sense of smell, not able of scenting’ \citep{Garnier2021}, eventually reanalyzed as a privative \textit{bahuvrīhi} \is{Bahuvrīhi} (viz.\ ‘lacking [sound] \iwi{AG}{φρένες}’), which implies that Hom.\ \iwi{AG}{φρένες}  would be a \isi{backformation}, coined after other names of body parts showing the suffix -ήν. In Homer, the \iwi{AG}{φρένες}  (this organ is considered double) are rather the seat of the intelligence, memory, and perception. When someone is angry, his \iwi{AG}{φρένες}  turn black, because they are \textit{filled} with anger, as in A 103, μένεoς  δὲ μέγα φρένες  ἀμφὶ μέλαιναι \# πίμπλαντ(o) “and with rage was his black heart wholly filled” (transl.\ Loeb).} The PIE root *\textit{bʰerh\textsubscript{2}}-\iw{PIE}{*\textit{bʰerh\textsubscript{2}}-/*\textit{bʰreh\textsubscript{2}}-} ‘sich schnell bewegen’ (\citetalias[ 81]{Rix2001})\footnote{cf.\ Ved.\ \iwi{Sanskrit}{\textit{bhar\textsuperscript{i}}-} ‘to move rapidly, to rush’, Hitt.\ \iwi{Hitt}{\textit{parḫ}-} ‘to chase, pursue, to hunt’, and the newly identified HLuw.\ */par.xa-/\iw{HLuw}{*\textit{parxa}-} ‘to expel’ (\citealt[204--206]{Melchert2016}).} is perhaps a better candidate for a noun ‘fury, anger’.\footnote{Ved.\ \textit{bhū́r-ṇi}- \iw{Sanskrit}{\textit{bhū́rṇi}-} [adj.] ‘restless, active, angry’.} As I have previously suggested (\citealt[53, fn.\ 6]{Garnier2021}), the PIE root \iwi{PIE}{*\textit{bʰerh\textsubscript{2}}-/*\textit{bʰreh\textsubscript{2}}-} ‘to hurry, to rush’ (with \textit{Schwebeablaut}) \is{Schwebeablaut} could be explained as a reanalyzed nominal stem \iwi{PIE}{*\textit{bʰér-h\textsubscript{2}}-/*\textit{bʰr-éh\textsubscript{2}}-} ‘hurry, haste, fury’ [of \isi{proterokinetic} inflection], built on the PIE primary root\is{derivation!deradical}  \iwi{PIE}{*\textit{bʰer}-} (cf.\ Gk.\ \iwi{AG}{φέρoμαι} [mid.] ‘to move rapidly’). 

\subsection{PIE *\textit{bʰr̥h\textsubscript{2}}\textit{-én-o}- [adj.] ‘being in quick motion’}
\label{sec:2.3}

Ved.\ *\textit{bhur-áṇ-a}- \iw{Sanskrit}{*\textit{bhuráṇa}-} [adj.] ‘quick, active’, similar to Ved.\ \textit{tur-áṇa-} \iw{Sanskrit}{\textit{turáṇa}-} [adj.] ‘swift’ for its formation, may represent something inherited, if one compares W.\ \iwi{Welsh}{\textit{bar-an}} [m.] ‘fury, anger’ (< Proto-Celt.\ \iwi{PCelt}{*\textit{barano}-}), which shows a trivial \isi{conversion} from *‘angry’ to ‘anger’ (perhaps via an old neuter stem \*\textit{baranon}). For the semantic discrepancy between ‘active’ and ‘angry’, one may mention Ved.\ \textit{bhū́r-ṇi}- \iw{Sanskrit}{\textit{bhū́rṇi}-} [adj.] ‘restless, active, angry’ (< PIE \iwi{PIE}{*\textit{bʰr̥h\textsubscript{2}-n-í}-} [f.] ‘vehemence’).

\subsection{PIE *\textit{bʰr̥h\textsubscript{2}}\textit{-én-t}- [adj.] ‘being in quick motion’}
\label{sec:2.4}

From PIE \iwi{PIE}{*\textit{bʰr̥h\textsubscript{2}-én}} [loc.] could be derived an adjective \iwi{PIE}{*\textit{bʰr̥h\textsubscript{2}-én-t}-} ‘being in quick motion’, with the “characterizing” or “individualizing” morpheme *-\textit{t}-,\iw{PIE}{*-\textit{(e)t}-} a delocative combination\is{delocatival} producing \iwi{PIE}{*-\textit{én-t}-} which could eventually be reanalyzed as a \isi{participle} *-\textit{ént}-\iw{PIE}{*-\textit{(e/o)nt}-} to the zero grade-stem of a root\is{root!present} present.\footnote{This derivational pattern is exemplified by PIE \iwi{PIE}{*\textit{h\textsubscript{1}d-én}} [loc.] ‘beim Beißen befindlich’ → hypostatic \iwi{PIE}{*\textit{h\textsubscript{1}d-én-t}-} ‘beißend’ (\citealt[150--151, fn.\ 211]{Neri2017}). One may also add several Vedic “aorist” participles\is{participle} not referring to a previous action, such as Ved.\ \textit{dhr̥ṣ-ánt}- \iw{Sanskrit}{\textit{dhr̥ṣánt}-} ‘bold’ \citep[163]{Lowe2017}, which could reflect PIE \iwi{PIE}{*\textit{dʰr̥s-én}} [loc.] ‘with boldness’ → \iwi{PIE}{*\textit{dʰr̥s-én-t}-} ‘bold’.} This very form is reflected by YAv.\ \iwi{Avestan}{\textit{barǝṇt}-} [ptcp.]\is{participle} ‘stürmend’ (= Ved.\ \iwi{Sanskrit}{*\textit{bhur-ánt}-} ‘quivering’ replaced by the thematic mid.\ ptcp.\ \is{participle} \iw{Sanskrit}{\textit{bhurámāṇa}-}  \textit{bhur-á-māṇa}-) and Hitt.\ \iwi{Hitt}{\textit{parḫ-ant}-}  [ptcp.]\is{participle} (< PA \iwi{PAn}{*\textit{barH-ánt}-}) meaning both: (i) ‘*being in quick motion, galloping’ [\isi{intransitive}]; (ii) ‘set in quick motion’ [\isi{transitive} passive]. The \isi{intransitive} meaning (i) is secure in: \textit{pár-ḫa-an-d}[\textit{a-an} [acc.sg.c.] ‘galloping’ (KUB 35.145 rev.\ 13), by contrast with (ii) \textit{arḫa pár-ḫa-an-za} [nom.sg.c.] ‘expelled/banished’ (KUB 8.1 ii 7–8). The scalar process of grammaticalization leading to a “pure” \isi{transitive} passive \isi{participle} (‘expelled’) developed by increasing the verb’s valency by turning it into a \isi{causative} (note the use of \iwi{Hitt}{\textit{arḫa}} [adv.] ‘off, away’ here with the lexical value of a \isi{preverb}). As a result, one may speculate that the \isi{transitive} present \textit{parḫ-zi} \iw{Hitt}{\textit{parḫ}-} ‘to chase’ is in fact backformed\is{backformation} after its passive \isi{participle} \iwi{Hitt}{\textit{parḫ-ant}-}* ‘chased’.

\subsection{PIE *\textit{bʰr̥h\textsubscript{2}-en-i̯ó-} [adj.] ‘being in quick motion’ and related forms}
\label{sec:2.5}

As already mentioned in Section \ref{sec:sota}, one has to posit a delocatival adjective\is{delocatival} \iwi{PIE}{*\textit{bʰr̥h\textsubscript{2}-en-i̯ó}-} [adj.] ‘being in quick motion’, indirectly reflected by its “characterized” \isi{doublet} \iwi{PIE}{*\textit{bʰr̥h\textsubscript{2}-en-i̯ó-nt}-} [adj.] ‘the swift one’. This very form was eventually reinterpreted as a present “active” \isi{participle} (Ved.\ \textit{bhur-aṇ-yánt}- \iw{Sanskrit}{\textit{bhuraṇyánt}-} ‘quivering, active’ [always of Agni]) of an \isi{intransitive} verb \textit{bhur-aṇ-yá-\textsuperscript{ti}} \iw{Sanskrit}{\textit{bhuraṇyá}-} ‘to be active or restless’. Independently, there was also a renewed adjective (Ved.\ \textit{bhur-aṇ-yú}- \iw{Sanskrit}{\textit{bhuraṇyú}-} [adj.] ‘quivering, active’ < PIE \iwi{PIE}{*\textit{bʰr̥h\textsubscript{2}-en-i̯ú}-}). The main bulk of the attestations consists of \isi{intransitive} forms: by contrast with the labile\is{labile verbs} \isi{transitive} present \textit{bhur-aṇ-yá-\textsuperscript{ti}} \iw{Sanskrit}{\textit{bhuraṇyá}-} ‘to agitate, stir (a liquid)’ [RV 2x], there is \textit{bhur-aṇ-yá-\textsuperscript{ti}} \iw{Sanskrit}{\textit{bhuraṇyá}-} ‘to be active or restless’ [RV 3x] (surely backformed\is{backformation} after  \iw{Sanskrit}{\textit{bhuraṇyánt}-}  \textit{bhuraṇ-yánt}-), \textit{bhur-aṇ-yánt}- ‘quivering, active’ [RV 2x], and \textit{bhur-aṇ-yú}- \iw{Sanskrit}{\textit{bhuraṇyú}-} [adj.] ‘quivering, active’ [RV 5x]. In other words, 10 out of 12 attestations are \isi{intransitive}.

\section{PIE 1.*\textit{bʰeh\textsubscript{2}}\textit{-} ‘to shine’}
\label{sec:3}

One may also reconstruct PIE \iwi{PIE}{*\textit{bʰh\textsubscript{2}-én}} [loc.] ‘into light; showing; appearing’, which was an adverb, perhaps compositional only (viz.\ PIE *°\textit{bʰh\textsubscript{2}-én}). This adverbial locative was presumably  built on a root noun\is{root!noun} ([virtual] PIE \iwi{PIE}{*\textit{bʰóh\textsubscript{2}}-} [f.] ‘light’), which is well-attested in composition, as is clear from Ved.\ \iwi{Sanskrit}{\textit{vi-bh{ā́}}-} [adj.] ‘shining’ (< PIE \iwi{PIE}{*\textit{h\textsubscript{1}u̯i-bʰéh\textsubscript{2}}-}). This -\textit{én}-locative \iw{PIE}{*-\textit{én}} would be the derivational base of an adjective *-\textit{bʰh\textsubscript{2}-én-t}-\iw{PIE}{*\textit{bʰh\textsubscript{2}-ént}-, *-\textit{bʰh\textsubscript{2}-én-t}-} ‘shining’ with the aforementioned “characterizing” or “individualizing” morpheme *-\textit{t}-\iw{PIE}{*-\textit{(e)t}-} (cf.\ Section \ref{sec:2.4}), for instance PIE \iwi{PIE}{*\textit{pro-bʰh\textsubscript{2}-én-t}-}  [adj.] ‘that appears before, manifest’ (derived from an adverb \iwi{PIE}{*\textit{pro-bʰh\textsubscript{2}-én}}), which was eventually reanalyzed as the “\isi{participle}” of an \isi{amphikinetic} present\is{root!present}  (PIE *\textit{bʰéh\textsubscript{2}-ti}, *\textit{bʰh\textsubscript{2}-énti} \iw{PIE}{1.\  *\textit{bʰéh\textsubscript{2}}-/\textit{bʰh\textsubscript{2}}-} ‘to shine’). One may assume that PIE \iwi{PIE}{*\textit{pro-bʰh\textsubscript{2}-én-t}-} [adj.] ‘shining forth’ → *\textit{pro-bʰh\textsubscript{2}-ént}- [ptcp.]\is{participle} would have given rise to a simplex *\textit{bʰh\textsubscript{2}-ént}- [ptcp.]\is{participle} ‘shining’. This form is reflected in the Hesychian gloss \iwi{AG}{φάντα}· \iwi{AG}{λάμπoντα} [acc.sg.m.] ‘shining’ (Gk.\ \iwi{AG}{φάντα} < PIE *\textit{bʰh\textsubscript{2}-ént-m̥}). It is noteworthy that this isolated form does not belong to any synchronic verbal stem: The etymologically related verb is φαίνω \iw{AG}{φαίνω, -oμαι} ‘to make appear, to show’ (φαίνoμαι [mid.] ‘to appear’). Within Greek, the privative adjective ἄ-φαν-τoς \iw{AG}{ἄφαντoς} ‘invisible’ and its sigmatic counterpart ἀ-φαν-ής \iw{AG}{ἀφανής} [adj.] ‘id.’ are deverbative formations\is{deverbal/deverbative!adjective} derived from the present stem φαίνω\iw{AG}{φαίνω, -oμαι} (< *\textit{pʰán-yō}).\iw{PG}{*\textit{pʰán-ye/o}-} Schirmer (\citetalias[ 69]{Rix2001}), after \citet[113]{Klingenschmitt1982}, assumes here a nasal-infixed present \is{nasal present} \iwi{PG}{*\textit{pʰán-e/o}-} (< PIE \iwi{PIE}{*\textit{bʰ-n̥-h\textsubscript{2}}-}) with \isi{thematization}, recharacterized with the suffix \iwi{PG}{*-\textit{ye/o}-}. This assumption would imply that Proto-Gk.\ \iwi{PG}{*\textit{pʰán-e/o}-} ‘to make visible’ would be closely related for its formation to Arm.\ \iwi{Armenian}{\textit{banam}} ‘to open’ (allegedly from PIE *\textit{bʰ-n-éh\textsubscript{2}-ti}, *\textit{bʰ-n̥-h\textsubscript{2}-énti}), \iw{PIE}{1.\ *\textit{bʰ-n-éh\textsubscript{2}}-/*\textit{bʰ-n̥-h\textsubscript{2}}-} and Albanian (Tosk \iwi[(Tosk)]{Albanian}{\textit{bëj}}, Gegh \iwi[(Gegh)]{Albanian}{\textit{bãj}}, Arbëresh, Arvanite \iwi[(Arbëresh, Arvanite)]{Albanian}{\textit{bënj}}) ‘to make appear; to do, to make’\footnote{For a summary of the attested forms see \citet{DemirajNeri2021}.} which points to Proto-Alb.\ \iwi{PAlb}{*\textit{ban-ja}-}.\footnote{By contrast with Vedic -\textit{anyá}-presents, \iw{Sanskrit}{-\textit{anyá}-} the majority of Proto-Alb.\ *-\textit{anja}-stems \iw{PAlb}{*-\textit{anja}-}  are \isi{transitive}, as ascertained by \citet[169--170]{Genesin2005} and \citet[50--51]{SchumacherMatzinger2013}.} However, Arm.\ \iwi{Armenian}{\textit{banam}} could be a productive\is{productivity} \iwi{Armenian}{-\textit{na}-}present of the type seen in Arm.\ \iwi{Armenian}{\textit{baṙnam}} ‘to lift’ (< \iwi{PArm}{*\textit{barj-nam}}) and \iwi{Armenian}{\textit{daṙnam}} ‘to come back’ (< \iwi{PArm}{*\textit{darj-nam}}), whereas Proto-Alb.\ \iwi{PAlb}{*\textit{ban-ja}-} ‘to bring to light; to produce’ could be derived from its \isi{participle} \iwi{PAlb}{*\textit{bana}-} (T.\ \iwi[(Tosk)]{Albanian}{\textit{bërë}}, G.\ \iwi[(Gegh)]{Albanian}{\textit{bënë}}), reflecting PIE \iwi{PIE}{*\textit{bʰh\textsubscript{2}-nó}-} [adj.] ‘shining, appearing’ or – alternatively – PIE \iwi{PIE}{*\textit{bʰh\textsubscript{2}-én-o}-} [adj.] ‘id.’). Besides, \iwi{PAlb}{*-\textit{nja}-}presents are productive\is{productivity} in Albanian and should therefore be rejected in comparisons. 

As a result, I will tentatively suggest another derivational chain for Gk.\ φαίνω\iw{AG}{φαίνω, -oμαι} ‘to make appear’. In my opinion, one may start with an adverb \iwi{PG}{*\textit{pró-pʰa}} ‘appearing before’ (< PIE \iwi{PIE}{*\textit{pro-bʰh\textsubscript{2}-én}})\footnote{Except for Hom.\ \iwi{AG}{*αἰϝέν} [adv.] ‘for ever’, which was oxytone, Greek has restructured the PIE locatives in *-\textit{en} \iw{PIE}{*-\textit{én}} on zero grade *-{\textit{n̥}}: see for instance Gk.\ \iwi{AG}{χεῖμα} [adv.] ‘in winter’ vs.\ Ved.\ \iwi{Sanskrit}{\textit{héman}} ‘id.’ (< PIE \iwi{PIE}{*\textit{g̑ʰéi̯-m-en}}), as \textit{per} \citet[189]{Nussbaum1986}. As a result, Gk.\ \iwi{PG}{*\textit{pró-pʰa}} [adv.] ‘appearing before’ (for expected ***\textit{pró-pʰan}) should be the reflex of PIE \iwi{PIE}{*\textit{pro-bʰh\textsubscript{2}-én}}. Interestingly, alongside \iw{PIE}{*-\textit{én}} *-\textit{én}, there is also evidence for \iw{PIE}{*-\textit{ér}} *-\textit{ér} used to build adverbs. Gk.\ \iwi{AG}{ἄφαρ} [adv.] ‘quickly’ could be explained by PIE \iwi{PIE}{*\textit{sm̥-bʰh\textsubscript{2}-ér}} [adv.] ‘in a glimpse’, an adverbial formation built on PIE \iwi{PIE}{*\textit{sm̥-bʰéh\textsubscript{2}}-} [f.] ‘glimpse’, rather than by PIE \iwi{PIE}{\textsuperscript{†}*\textit{h\textsubscript{2}ébʰ-r̥}} ‘rapid stream’ assumed by \citet[327]{Willi2004}.} as the base stem of a denominative verb\is{denominative!verb} \iwi{AG}{πρo-φαίνω} (< *\textit{pro-pʰán-yō})\iw{PG}{*\textit{pro-pʰán-ye/o}-} ‘to make appear before’, displaying the same pattern as Hom.\ \iwi{AG}{λίπα} [adv.] ‘richly, in a fat way’ (of anointing) → \iwi{AG}{λιπαίνω} ‘to anoint’. One may go a step further and assume that the “complex” suffix \iwi{PG}{*-\textit{nye/o}-}  is due to resegmentation of the present stem \iwi{PG}{*\textit{pro-pʰán-ye/o}-}  as *\textit{pro-pʰá-nye/o-}, after which were coined such presents as Gk.\ \iwi{AG}{κλῑ́νω} (< \iwi{PG}{*\textit{klí-nyō}})  ‘to cause to lean’ (with sigmatic deverbatives\is{deverbal/deverbative!adjective} in *°\textit{klĭn-ḗs}) and \iwi{AG}{πλῡ́νω} (< \iwi{PG}{*\textit{plú-nyō}}) ‘to wash, to clean’ (with sigmatic deverbatives\is{deverbal/deverbative!adjective} in *°\textit{plŭn-ḗs}).\footnote{In the final chapter of his thesis, \citet[121--135]{Praust1998} reconstructs an athematic -\textit{ni}-present class: for instance, PIE *\textit{\kinvbreve li-néi̯-ti}, *\textit{\kinvbreve li-ni̯-énti} \iw{PIE}{*\textit{\kinvbreve li-néi̯}-/*\textit{\kinvbreve li-ni̯}-} to account for Gk.\ \iwi{AG}{κλῑ́νω} (< \iwi{PG}{*\textit{klí-nyō}}) ‘to cause to lean’ (Aeol.\ \iwi[(Aeol.)]{AG}{κλίννω}). The adduced comparanda are YAv.\ \iwi[(YAv.)]{Avestan}{\textit{siri-nu}-} (or \textit{sri-nu}-) and Lat.\ \iwi{Latin}{°\textit{clīnāre}} (with a long -\textit{ī}-), which he considers to be a \isi{durative} present to a simplex \iwi{Latin}{*\textit{cliniō}} \citep[132]{Praust1998}. This is simply an \textit{ad hoc} assumption.}

A possible reflex of PIE *\textit{(pro-)bʰh\textsubscript{2}}\textit{-én} \iw{PIE}{*\textit{pro-bʰh\textsubscript{2}-én}} [adv.] may be attested in Germanic: Go.\ \iwi{Goth}{\textit{bandwjan}} ‘to give a sign’ and \iwi{Goth}{\textit{bandwa}*} [str.f.] ‘sign’ could be derived from a Proto-Germanic action noun \iwi{PGmc}{*\textit{bandu}-} [m.] ‘bright appearance’ (< PIE \iwi{PIE}{*\textit{bʰh\textsubscript{2}-en-tú}-} [abstract] ← \iwi{PIE}{*\textit{bʰh\textsubscript{2}-en-tó}-} [adj.] ‘endowed with shine’). \textit{Pace} \citetalias[ 51--52]{Kroonen2013}, Go.\ \iwi{Goth}{\textit{bandwjan}} and \iwi{Goth}{\textit{bandwa}*} can hardly be related to Proto-Gmc.\ \iwi{PGmc}{*\textit{bann-an}-} [str.v.] ‘to command, to summon’ (for which see Section \ref{sec:4}). The original meaning was obviously ‘apparition, vision’, as is clear from the Langobardic cognate of Go.\ \iwi{Goth}{\textit{bandwa}*} borrowed in Late Latin (\textit{vexillum quod bandum appellant} in Paulus Diaconus). The concrete meaning ‘flag, banner’ is also supported by Byz.\ Gk.\ \iwi[(Byz.)]{AG}{βάνδoν} [nt.] (in Procopius). For the meaning, one may compare Ved.\ \iwi{Sanskrit}{\textit{ketú}-} [m.] ‘bright appearance’ and ‘flag, banner’ (< PIE \iwi{PIE}{*\textit{k\textsuperscript{u̯}oi̯t-ú}-}).

\section{PIE 2.\ *\textit{bʰeh\textsubscript{2}}\textit{-} ‘to speak, to say’}
\label{sec:4}

From PIE \iwi{PIE}{2.\ *\textit{bʰeh\textsubscript{2}}-} ‘to speak, to say’, we may assume an adverbial oblique \iwi{PIE}{*\textit{bʰh\textsubscript{2}-én}} as the derivational base of a  PIE adjective  \iwi{PIE}{*\textit{bʰh\textsubscript{2}-én-o}-}  ‘speaking’ whose characterized form \iwi{PIE}{*\textit{bʰh\textsubscript{2}-én-o-nt}-} ‘the speaking one; the one endowed with speaking’\footnote{Likewise, after Hitt.\ \iwi{Hitt}{\textit{ḫappina}-} [adj.] ‘rich’ (< PIE \iwi{PIE}{*\textit{h\textsubscript{3}ép-en-o}-}) has been coined \iwi{Hitt}{\textit{ḫappinant}-} [c.] ‘rich person’ (< [virtual] PIE \iwi{PIE}{*\textit{h\textsubscript{3}ép-en-o-nt}-} ‘the rich one’).} is surely reflected by Ved.\ \textit{bhán-ant}-*\iw{Sanskrit}{\textit{bhánant}-*} [ptcp.]\is{participle} which would thus antedate the indicative \textit{bhán-a-ti} \iw{Sanskrit}{\textit{bhána}-} ‘to speak, to say’. This  stem has been accounted for by Schirmer (\citetalias[ 69--70 \& fn.\ 6]{Rix2001}) as the restructuring of an athematic nasal-infixed present \is{nasal present} *\textit{bhan{ā́}-ti}, *\textit{bhan-ánti} (< PIE 2.\ *\textit{bʰ-n-éh\textsubscript{2}-ti}, *\textit{bʰ-n̥-h\textsubscript{2}-énti} \iw{PIE}{2.\ *\textit{bʰ-n-éh\textsubscript{2}}-/*\textit{bʰ-n̥-h\textsubscript{2}}-} ‘to speak, to say’). Kroonen (\citetalias[ 52]{Kroonen2013}) instead assumes a thematic root present\is{thematic!present} \mbox{\iwi{PIE}{*\textit{bʰénh\textsubscript{2}-e/o}-},} which he compares to Proto-Gmc.\ \iwi{PGmc}{*\textit{bann-an}-} ‘to summon’, allegedly reflecting IE \textsuperscript{†}\textit{bʰónh\textsubscript{2}}\textit{-e}/\textit{o-}, with the final laryngeal triggering the gemination of the nasal. He does not mention the reconstruction of this strong verb as PIE  \iwi{PIE}{*\textit{bʰh\textsubscript{2}-nu̯-é/ó}-} by \citet[89]{Seebold1970}, which was perhaps less anachronistic than its rewriting as IE \textsuperscript{†}\textit{bʰe-nu̯-é/ó-} by Kroonen (\citetalias[ 52]{Kroonen2013}), who tacitly suggests a verb \iwi{PGmc}{*\textit{binn-an}-} ‘to speak’ endowed with an \isi{intensive} \isi{doublet} \iwi{PGmc}{*\textit{bann-an}-} ‘to summon’. As already mentioned (see Section \ref{sec:3}), the etymological connection with Go.\ \iwi{Goth}{\textit{bandwo}, \textit{bandwa}*} [f.] ‘sign’ has little to recommend itself, not to mention the etymon \textsuperscript{†}\textit{bʰonh\textsubscript{2}-tú-h\textsubscript{2}-} offered by the author.

I would rather assume here PIE \iwi{PIE}{*\textit{bʰh\textsubscript{2}-en-u̯ó}-} ‘endowed with speaking’, a decasuative adjective\is{decasuative} of the pattern *\textit{CC-en-u̯ó}- which is required by Proto-Gk.\ \iwi{AG}{*ξέν-ϝoς}  [m.] ‘guest’ < PIE \iwi{PIE}{*\textit{gʰs-en-u̯ó}-} [adj.] ‘being at the feast’ \citep[65]{Garnier2013}.\footnote{PIE \iwi{PIE}{*\textit{gʰs-én}} [loc.] ‘at the feast’ may be assigned to PIE \iwi{PIE}{*\textit{gʰés-ōr}} [AK \is{amphikinetic} collective] ‘rich meal’ derived from \iwi{PIE}{*\textit{gʰés-r̥}} [PK nt.] ‘action of eating’ \citep[65]{Garnier2013}. Note that Ir.\ \iwi{PIr}{2.\ *\textit{xšnaw}-} ‘to welcome a guest’ (\citealt[457--458]{Cheung2007}) is possibly backformed\is{backformation} after its ptcp.\ \is{participle} \iwi{PIr}{*\textit{xšnu-tá}-} ‘invited, welcome’ (< [virtual] PIE \iwi{PIE}{*\textit{gʰs-n-u-tó}-}). Perhaps one has to assume here PIE \iwi{PIE}{*\textit{gʰs-en-u̯-ó-nt}-} [characterized adj.] secondarily inflected according to the \isi{amphikinetic} inflection (viz.\ \is{genitive}gen.sg. *\textit{gʰs-n̥-u̯-n̥t-ós} resyllabified in *\textit{gʰs-nu-u̯-n̥t-ós} {>}{>} Ir.\ \iwi{PIr}{*\textit{xšnuw-ánt}-} ‘invited, welcome’ normalized as as passive ptcp.\ \is{participle} \iwi{PIr}{*\textit{xšnu-tá}-} ‘id.’). More recently, \citet[338]{Steer2019} has rejected the derivation \iwi{PIE}{*\textit{gʰs-en-u̯ó}-} [adj.] ‘being at the feast’, suggesting instead a PIE deverbative\is{deverbal/deverbative!adjective} \iwi{PIE}{*\textit{gʰs-nu̯-ó}-} [adj.] ‘eating’ to a -\textit{nu}-present \is{nasal present} assumed by \citet[284, 286]{Lipp2009} → \textit{vr̥ddhi}-derivative \is{\textit{vr̥ddhi}} (viz.\ PIE *\textit{gʰs-énu̯-o}- ‘eater’), which seems impossible to me: One would expect PIE *\textit{gʰs-neu̯-ó}-. See the discussion of Steer’s proposal by \citet[160--161]{Schmitt2020}, who considers that a base meaning like *‘eater’ is perhaps too trivial to account for the religious aspects of Homeric hospitality. Schmitt duly mentions the figure of Zεὺς  ξένιoς. I had addressed this issue by assuming that the base is PIE \iwi{PIE}{*\textit{gʰós}-} [f.] ‘rich meal ceremoniously offered to the guests’.} One may therefore assume a “characterized” \iwi{PIE}{*\textit{bʰh\textsubscript{2}-en-u̯ó-nt}-} ‘the one who speaks with authority; the one who gives orders’ surfacing as Post-IE *\textit{bʰa-nu̯-ónt}- [m.] ‘commanding [adj.], commander [m.]’. This form was normalized as Proto-Gmc.\ \iwi{PGmc}{*\textit{bánn-and}-} [ptcp.]\is{participle} with secondary root accent. This was the starting point for a strong verb \iwi{PGmc}{*\textit{bann-an}-} ‘to summon’. Such a verbal reinterpretation of an adjectival stem could have been fostered by other genuine *\textit{and}-participles\is{participle} \iw{PGmc}{*-\textit{and}-} used as agent nouns in Germanic: see for instance OE \iwi{OE}{\textit{būend}} [m.] ‘dweller, inhabitant’.

Incidentally, one may speculate that the Attic infinitive \iwi{AG}{φάναι} (vs.\ Hom.\ mid.\ inf.\ \iwi{AG}{φάσθαι}) may be the recharacterization of Proto-Gk.\ \iwi{PG}{*\textit{pʰán}} ‘in speaking’ (< PIE \iwi{PIE}{*\textit{bʰh\textsubscript{2}-én}}), like the infinitive \iwi{AG}{ἰέναι} ‘to go’ is the reflex of \iwi{PIE}{*\textit{ién}} ‘on the way’ enlarged with the Greek allative adverbial suffix \iwi[(allat.)]{AG}{-αι} (cf.\ Hom.\ \iwi{AG}{ἴμεν} vs.\ \iwi{AG}{ἴμεναι}), as suggested by  \citet[323]{GarnierPinault2020}.

\section{PIE *\textit{g\textsuperscript{u̯}}\textit{eh\textsubscript{2}}- ‘to make a step; to be stable’}

One may reconstruct PIE \iwi{PIE}{*\textit{g\textsuperscript{u̯}h\textsubscript{2}-én}} [adv.] ‘firmly standing’ as the base of \iwi{PIE}{*\textit{g\textsuperscript{u̯}h\textsubscript{2}-én-t}-} [adj.] ‘stable’ (reanalyzed as an “aorist” \isi{participle} *\textit{g\textsuperscript{u̯}h\textsubscript{2}-ént}-). This form is reflected by Hom.\ εὖ \iwi{AG}{δια-βᾱ́ς}  (M 458) ‘setting his feet well apart’, with a clear static value, and not referring to a prior event (\citealt[133--135]{Letoublon1986}; \citealt[935]{Lamberterie1990}). For the semantics, one may compare Gk.\ \iwi{AG}{βέβαιoς}  [adj.] ‘firm, stable’. There was possibly also a decasuative\is{decasuative} \iwi{PIE}{*\textit{g\textsuperscript{u̯}h\textsubscript{2}-en-i̯ó}-} [adj.] ‘firmly standing’ endowed with a “characterizing” \isi{doublet} \iwi{PIE}{*\textit{g\textsuperscript{u̯}h\textsubscript{2}-en-i̯ó-nt}-} ‘the one firmly standing’. This adjectival formation could account for Gk.\ (εἰσ-, ἐμ-, ἐπι-)\iwi{AG}{βαίνoντ-} [ptcp.]\is{participle} ‘stepping; setting his feet’ (Hom.\ πoσὶ βαίνων). One must not confuse Gk.\ \iwi{AG}{1.\ συμβαίνειν ‘to stand with the feet together’} (of a statue) with \iwi{AG}{2.\ συμβαίνειν ‘to happen’} (ultimately from PIE \iwi{PIE}{*\textit{g\textsuperscript{u̯}m̥-i̯é/ó}-} ‘to come [in contact with]’). Similarly, Gk.\ \iwi{AG}{σύμβασις}  [f.] ‘agreement’ indirectly reflects PIE \iwi{PIE}{*\textit{kóm-g\textsuperscript{u̯}m̥-ti}-} [f.] ‘coming together; agreement’, semantically close to Osc.\ \iwi{Oscan}{\textbf{kúmbened}} ‘it is agreed’ (Lat.\ \textit{convēnit}\iw{Latin}{\textit{conuēnit}}), but is not related to \iwi{AG}{βάσις}  [f.] ‘stepping, step, steps’, which rather reflects PIE \iwi{PIE}{*\textit{g\textsuperscript{u̯}h\textsubscript{2}-tí}-} [f.] ‘action of stepping’. Both roots have merged in Greek, due to the so-called “satellite-framing” (for instance πρo-βαίνω ‘to step forward, to advance’ is \textit{dynamic}). As a result, the routinely adduced equation between Gk.\ \iwi{AG}{°βαίνω} ‘to step, to walk’ and Lat.\ \textit{veniō}\iw{Latin}{\textit{ueniō}} ‘to come’ (< PIE \iwi{PIE}{*\textit{g\textsuperscript{u̯}m̥-i̯é/ó}-}), as \textit{per} Kümmel (\citetalias[ 209]{Rix2001}) is only partly true.

\section{Concluding remarks} 
\largerpage

The above study highlights the process of grammaticalization of “petrified” decasuative formations\is{decasuative} derived from *\textit{CC}-\textit{én} locatives of root nouns\is{root!noun} (especially in composition). By internal or external derivation, such forms could produce various delocatives:\is{delocatival} 

\begin{exe}
    \ex \begin{xlist} \ex  PIE *\textit{CC-ḗn}  \iw{PIE}{*-\textit{ḗn}} [HK f.] (< PIE *\textit{CC-én-s}) \label{ex:1a}
\ex  PIE *\textit{CC-én-t-} \iw{PIE}{*-\textit{én-t}-} [adj.]\\\null \hspace{0.5cm} (\textbf{→} pseudo-ptcp.\ \is{participle} *\textit{CC-ént-} \iw{PIE}{*-\textit{(e/o)nt}-}   [root present\is{root!present} or root\is{root!aorist} aorist])
\ex PIE *\textit{CC-én-o-} \iw{PIE}{*-\textit{én-o}-} [adj.]\\\null   \hspace{0.5cm} \textbf{→} *\textit{CC-én-o-nt-} \iw{PIE}{*-\textit{én-o-nt}-} [characterized adjective]\\ \null \hspace{1cm} (\textbf{→} pseudo-ptcp.\ \is{participle} *\textit{CCén-ont-})
\ex PIE *\textit{CC-en-i̯ó-} \iw{PIE}{*-\textit{en-i̯ó}-} [adj.]\\\null  \hspace{0.5cm} \textbf{→} *\textit{CC-en-i̯ó-nt-} \iw{PIE}{*-\textit{en-i̯ó-nt}-} [characterized adjective]\\ \null \hspace{1cm} (\textbf{→} pseudo-ptcp.\ \is{participle} *\textit{CCen-i̯ónt-})
\ex PIE *\textit{CC-en-u̯ó-} \iw{PIE}{*-\textit{-en-u̯ó}-}  [adj.]\\\null  \hspace{0.5cm} \textbf{→} *\textit{CC-en-u̯ó-nt-} \iw{PIE}{*-\textit{en-u̯ó-nt}-} [characterized adjective]\\ \null  \hspace{1cm} (\textbf{→} pseudo-ptcp.\ \is{participle} *\textit{CCénu̯-ont-}) \label{ex:1e}
\ex PIE *\textit{CC-en-tó-} \iw{PIE}{*-\textit{en-tó}-}  [adj.]\\ \null  \hspace{0.5cm} \textbf{→} abstract *\textit{CC-en-tú}- \iw{PIE}{*-\textit{en-tú}-} [m.]
\end{xlist}
\end{exe}



Only the patterns (\ref{ex:1a}--\ref{ex:1e}) could produce -\textit{nt}-stems\iw{PIE}{*-\textit{(e/o)nt}-}  which were eventually reanalyzed as participles\is{participle} of finite verbs. This new morphological analysis may account for Ved.\ \textit{bhur-aṇ-yánt-} \iw{Sanskrit}{\textit{bhuraṇyánt}-}  ‘quivering, active’, Gk.\ φαίνω\iw{AG}{φαίνω, -oμαι} ‘to make appear’, \iwi{AG}{°βαίνω} ‘to step, to walk’, and Proto-Gmc.\ \iwi{PGmc}{*\textit{bann-an}-} ‘to summon’, which were previously lacking a clear derivational analysis.

\section*{Abbreviations}

\begin{tabularx}{.45\textwidth}{lQ}
adj.&  adjective \\
 adv. & adverb \\
 Alb.\ & Albanian \\
 Arm.\ & Armenian \\
 Av.\ & Avestan \\
Byz. Gk.\ & Byzantine Greek \\
 c. & common gender \\
 CLuw.\ & Cuneiform Luwian \\
 f. & feminine \\
 Gk.\ & Greek\\
 Go.\ & Gothic \\
 Hitt.\ & Hittite \\
 HK & hysterokinetic \\
 Hom.\ & Homeric \\
 Lat.\ & Latin \\
 loc. & locative \\
 \end{tabularx}
\begin{tabularx}{.5\textwidth}{lQ}
 Lyc.\ & Lycian\\
  m. & masculine\\
  nt. & neuter\\
  Osc. & Oscan\\
  OE & Old English\\
  OIr.\ & Old Irish\\
 PA & Proto-Anatolian\\
  PIE & Proto-Indo-European\\
   Pre-Lyc.\ & Pre-Lycian\\
  Proto-Alb.\ & Proto-Albanian\\
  Proto-Celt.\ & Proto-Celtic\\
   Proto-Gk.\ & Proto-Greek\\
  Proto-Gmc.\ & Proto-Germanic\\
 Ved.\ & Vedic\\
 YAv.\ & Young Avestan\\
 W. & Welsh\\
 \end{tabularx}




\sloppy\printbibliography[heading=subbibliography,notkeyword=this]
\end{document} 
 









